\documentclass[../main]{subfiles}
\begin{document}
\chapter{Principio de Arquímedes}
\img{images/02/arquimedes}{Buque de transporte.}
\info{
Contenido}{
Principio de Arquímides.
}

\info{Principio de Arquímides}{El principio de Arquímedes afirma que todo cuerpo sumergido en un fluido experimenta un empuje vertical y hacia arriba igual al peso de fluido desalojado.
\begin{equation}
    E=\gamma \times V
\end{equation}

Siendo:

$E: $Empuje [kg]

$\gamma:$ Peso específico del líquido$[\vec{kg}/m^3]$

$V:$ Volumen $[m^3]$
}

\info{Peso aparente}{
 El peso aparente $P_a$ es igual al peso real que tiene el cuerpo fuera del fluído menos el peso del fluido que desplaza al sumergirse.
\begin{equation}
    P_a=P-E= V \times \gamma_{cuerpo} - V \times \gamma_{liquido}= V (\gamma_{cuerpo}-\gamma_{liquido})
\end{equation}
Siendo:

$P_a:$ Peso aparente.

$P$ : Peso Real, fuera del líquido.

$E$ : Peso del líquido desalojado.


}

\actividad{Resolver los siguientes problemas.}
\begin{enumerate}
% de internet https://www.fisimat.com.mx/principio-de-arquimedes/
\item Un cubo de hierro de 20cm de arista se sumerge totalmente en agua $\gamma_{H_2O}=1000\vec{kg}/m^3$. Si tiene un peso con una magnitud de 560.40 N, calcular: ¿Qué magnitud de empuje recibe y cuál será la magnitud del peso aparente del cubo?

\item Una esfera de volumen de $V_e=3 \times 10^{-4}m^3$, está totalmente inmersa en un líquido cuya densidad es de $\rho_{l}=900 kg/m^3$, determine, a) La intensidad de empuje que actúa en la esfera, b) La intensidad del peso de la esfera para que se desplaza hacia arriba o hacia abajo. 

\item Un cubo de cobre, de base igual a $b=35cm^2$ y una altura de $h=12 cm$, se sumerge hasta la mitad, por medio de un alambre, en un recipiente que contiene alcohol.  ¿Cuál es la magnitud del peso aparente del cubo debido al empuje, si la magnitud de su peso es de 32.36 N?

\item Hay un iceberg flotando en el agua del mar $\rho_{mar}=1025 Kg/m^3$ de $V=60m^3$ del cual 2/3 esta sumergido. Calcular la masa del iceberg.

\img{./images/02/iceberg}{fotografía del iceberg que hundió el Titanic }

\item Durante un experimento, un cubo de madera de arista de $l=1m$ , se coloca en un recipiente que contiene agua. Se notó que el cubo flotó con el 60\% de su volumen sumergido. a) Calcule la intensidad del empuje ejercido por el agua sobre el bloque de madera, b) Calcule la intensidad de fuerza vertical $F$, que debe actuar sobre el bloque, para que permanezca totalmente sumergido.  

% otro de internet https://www.profesor10demates.com/2015/01/principio-de-arquimedes-ejercicios.html





\item Una esfera de $r=0,3m$ de radio flota en un recipiente con aceite ($\rho_{aceite}=800kg/m^3$). Si la esfera está sumergida hasta la mitad, calcular el peso de la misma.

\item Un objeto tiene una masa de $m=10Kg$ y ocupa un volumen de $V=7 litros$ , tiene un peso aparente de $P_a=24N$ dentro del líquido . Calcula la densidad del líquido 

\item ¿Sabrías decir cual es el peso aparente de un cubo de $10cm$ de lado y $10kg$ de masa que se sumerge completamente en un fluido cuya densidad es  $\rho=1000kg/m^3$?

\item  Calcula el volumen que se encuentra sumergida una piragua, si esta tiene una masa de $m=65Kg$ si la densidad del agua es $\rho=1000 kg/m^3$).
% inventados propios
\item El volumen pulmonar total de una persona es de $V=5,8lts$, si tiene un peso de $P=68\vec{kg}$ y ocupa un volumen de $V=66,4lts$ con los pulmones vacíos, averiguar si flota o se hunde con los pulmones llenos de aire y vacío.

\img{./images/02/pelps.jpg}{Michael Phelps nadador olímpico.}

    \item Calcular la profundidad de inmersión de un témpano de hielo de forma de paralelepípedo el cual emerge 5m sobre el nivel del mar, si el peso específico del hielo es $\gamma_h=975\frac{kg}{m^3}$ y el peso específico del agua del mar es de $\gamma_a=1,025\frac{kg}{m^3}.$
    
    \item Al sumergir un cuerpo que pesa $15\vec{g}$ en un líquido de peso especifico $0,9\frac{\vec{g}}{cm^3}$, se observa una pérdida de peso de $5\vec{g}$. Calcular el peso específico $\gamma$ del cuerpo.
    \item Un cuerpo pesa $P=15\vec{kg}$ reduce su peso a $P_{agua}=12\vec{kg}$ en el agua y $P_{liquido}=12,6\vec{kg}$ en otro cierto líquido. Calcular el peso específico. a) del líquido $\gamma_{liquido}$;b) del cuerpo $\gamma_{cuerpo}$. 
    \item Una cadena de oro y plata pesa $P=50\vec{g}$, en el agua reduce su peso a $P_{agua}=47\vec{g}$. Calcular la cantidad de oro y plata que posee si el peso específico del oro es $\gamma_{Au}=19,33\frac{\vec{g}}{cm^3}$ y el de la plata $\gamma_{Ag}=10,5\frac{\vec{g}}{cm^3}$.
    \item Un recipiente de 1 litro de capacidad se llena mediante dos líquidos $A$ y $B$, que en conjunto tienen un peso específico de $\gamma=1,2\frac{\vec{g}}{cm^3}$. El peso específico de $A$ es $\gamma_A=0,6\frac{\vec{g}}{cm^3}$ y el de $B$ es $\gamma_B=1,5\frac{\vec{g}}{cm^3}$.Calcular el volumen que ocupa cada componente.
    
    
    \item Una esfera de corcho de peso específico de $\gamma_{corcho}=0,20\frac{\vec{kg}}{dm^3}$ se sumerge en una profundidad de $h=10m$ de agua de peso específico $\gamma_{agua}=1\frac{\vec{kg}}{dm^3}$. Calcular el tiempo que se demora en llegar a la superficie, suponiendo que no existe resistencia al deslizamiento y existe una aceleración normal gravitatoria. 
    
    \img{./images/02/bolacorcho}{Esfera de corcho.}
    
    \item Un vaso vacío de $15cm$ de alto se sumerge verticalmente, con su boca hacia abajo, en una pileta llena de agua, de peso específico $\gamma_{H_2O}=1000\frac{\vec{kg}}{m^3}$, hasta que su fondo toca la superficie libre del agua. Calcular la altura $x$ de la columna de agua que penetra dentro del vaso. La presión atmosférica es de $P=1\frac{\vec{kg}}{cm^2}$.
    

\end{enumerate}

\end{document}